\scanpage{68}
% loc: 1–5, top, including interlinear additions
Ainsi, \deleted{\unsure{les}} \inserted{chaque} classe $C_i^*\deleted{\unsure{(e)}}$ \inserted{$(=C^\times(e))$} est munie \deleted{\unsure{chacune}} \deleted{(pour $e$ fixé)} d'une structure \unsure{triangulée|triangul.|triangulaire} \deleted{\unsure{non}} orientée canonique, \inserted{(indépendante du choix de $e$)} avec en plus une $\lambda$-orientat. \inserted{(dépendant à priori du choix de $e$)} qu'il faudrait expliciter. Le groupe projectif de $M$

% loc: 6
\[
 M^*/k^\times\quad(\simeq\operatorname{End}(E)/k^\times)
\]

% loc: left, oblique note beside 6–12
% slanted
\margin{c'est une pure \unsure{pétit\textsuperscript{o}|pét.|preuve}, il faut voir que si $u,v\in C^*$ sont \deleted{\unsure{tels}} tels que $\operatorname{Tr}uv=-1$, ils appartiennent à une $=$ classe $\ldots$ \unsure{cf|c'est} plus bas $\ldots$}

% loc: 7–8
permute entre \deleted{\unsure{elles|deux}} eux les $C_i^*\deleted{\unsure{(e)}}$, en respectant \inserted{(par transport des structures)} les structures orientées.

% loc: 9–13
On fera attention que pour \deleted{les autom qui \unsure{proviennent de}} $u\in\mathrm{Gl}(E)$, définissant $\bar u\in\mathrm{Gl}(E)/k^\times\simeq M^\times/k^\times$, l'endom. de $\deleted{\unsure{\mathrm{Gl}^{\times}(E)|\mathrm{Gl}(E)}}M^*$ \deleted{\unsure{qu}} (donc de $C^\times$) qu'il définit, est donné par
\[
 \theta\longmapsto u(\theta)=u\theta u^{-1}
\]

% loc: left at 14–18, double vertical rule
% slanted
\margin{\uline{Attention}}

% loc: 14–18
Ce n'est compatible avec l'application
\[
 \varphi^e:E^*\xrightarrow{\sim} C^\times(e)\subset C^*
\]
(\unsure{avec|car} $\mathrm{Gl}(E)$ opérant sur $E^\times$ et sur $C^\times$) que sur le sous-groupe $\mathrm{Sl}(E)$. De façon précise, on aura

% loc: 19, boxed
\[
 \boxed{\theta(\varphi^e(x))=\det\theta\,(\varphi^e(\theta(x)))}
\]

% loc: 20–21, bottom
\[
 \bigl(\ \theta(x\overset{e}{\otimes}x)=
 \theta(x)\overset{\theta(e)}{\otimes}\theta(x)=
 \theta(x)\overset{\det\theta.e}{\unsure{\odot|\otimes}}\theta(x)
 =\det\theta\bigl(\theta(x)\overset{e}{\otimes}\theta(x)\bigr)
\]
\deleted{Donc le groupe}
% archive sheet number: 67

